% ============================================================
\subsection{Architecture du Modèle Auto-Supervisé}
\label{sub:method}
% ============================================================
 
\subsubsection{Vue d'ensemble}
 
Nous proposons une adaptation unidimensionnelle du cadre
\emph{Image-based Joint-Embedding Predictive Architecture}
(I-JEPA, \cite{assran2023self}) à l'analyse statique de
binaires ELF. La représentation d'entrée est une \emph{séquence
de tokens VEX/API} de longueur fixe $L$, notée
\[
  \bm{x} = (x_1, x_2, \ldots, x_L) \in \mathcal{V}^{L},
\]
où $\mathcal{V}$ est le vocabulaire discret de taille $|\mathcal{V}|
= 371$ (voir Section~\ref{tab:dataset_stats}).
Le modèle comprend trois composants appris : un \emph{encodeur de
contexte} $f_\theta$, un \emph{encodeur cible} $f_\xi$ (également
appelé \emph{professeur}), et un \emph{prédicteur} $g_\phi$.
L'architecture globale est illustrée en Figure~\ref{fig:ijepa_overview}.
 
% ------------------------------------------------------------------
% Figure 1 : vue d'ensemble I-JEPA 1D  — layout vertical, espacé
% ------------------------------------------------------------------
\begin{figure}[ht]
\centering
\begin{tikzpicture}[
  every node/.style={font=\small},
  %% boîtes de modules
  enc/.style={
    draw, rounded corners=5pt,
    minimum width=3.4cm, minimum height=1.4cm,
    align=center, line width=1pt, text width=3.0cm
  },
  pred/.style={
    draw, rounded corners=5pt,
    minimum width=3.0cm, minimum height=1.2cm,
    align=center, line width=1pt, text width=2.6cm
  },
  repr/.style={
    draw, rounded corners=3pt,
    minimum width=3.2cm, minimum height=0.65cm,
    align=center, font=\small, line width=0.7pt
  },
  loss/.style={
    draw, rounded corners=5pt,
    minimum width=6.2cm, minimum height=0.8cm,
    align=center, line width=1pt, text width=5.8cm
  },
  %% cellules de séquence
  tok/.style={
    draw=black!35, minimum width=0.58cm, minimum height=0.58cm,
    inner sep=0pt, font=\scriptsize, rounded corners=1pt
  },
  %% flèches
  arr/.style={-Stealth, line width=1.0pt},
  emaArr/.style={-Stealth, line width=1.0pt, dashed, colTgt},
  sgLbl/.style={font=\scriptsize, colTgt!80!black, rotate=90},
]
 
% ═══════════════════════════════════════════════════════════════
% RANGÉE 1 — Séquences d'entrée (côte à côte, bien espacées)
% ═══════════════════════════════════════════════════════════════
 
% -- Colonne gauche : séquence masquée --------------------------
\node[font=\footnotesize\bfseries, anchor=center]
  (lbl_left) at (-3.4, 7.2) {Séquence $\tilde{\bm{x}}$ (masquée)};
 
% tokens : contexte (bleu), pré-cible (rouge clair), cible (rouge),
%          padding (gris)
\foreach \i/\col/\lbl in {
  0/colCtx!25/{\scriptsize VEX},
  1/colCtx!25/{\scriptsize OP},
  2/colMask!30/{\scriptsize[M]},
  3/colMask!30/{\scriptsize[M]},
  4/colMask!65/{\scriptsize[M]},
  5/colMask!65/{\scriptsize[M]},
  6/colPad!20/{\scriptsize[P]},
  7/colPad!20/{\scriptsize[P]}}{
    \node[tok, fill=\col] (tL\i) at (-5.1 + \i*0.70, 6.6) {\lbl};
}
 
% accolades sous tokens gauche
\draw[decorate,decoration={brace,mirror,amplitude=4pt}]
  (tL0.south west) -- (tL1.south east)
  node[midway,below=4pt,font=\tiny,colCtx!70!black]{contexte};
\draw[decorate,decoration={brace,mirror,amplitude=4pt}]
  (tL2.south west) -- (tL3.south east)
  node[midway,below=4pt,font=\tiny,colMask!70!black]{pré-cible};
\draw[decorate,decoration={brace,mirror,amplitude=4pt}]
  (tL4.south west) -- (tL5.south east)
  node[midway,below=4pt,font=\tiny,colMask!80!black]{cible};
\draw[decorate,decoration={brace,mirror,amplitude=4pt}]
  (tL6.south west) -- (tL7.south east)
  node[midway,below=4pt,font=\tiny,colPad!60!black]{pad};
 
% -- Colonne droite : séquence non masquée ----------------------
\node[font=\footnotesize\bfseries, anchor=center]
  (lbl_right) at (3.4, 7.2) {Séquence $\bm{x}$ (non masquée)};
 
\foreach \i/\col/\lbl in {
  0/colCtx!25/{\scriptsize VEX},
  1/colCtx!25/{\scriptsize OP},
  2/white/{\scriptsize API},
  3/white/{\scriptsize JK},
  4/colTgt!30/{\scriptsize VEX},
  5/colTgt!30/{\scriptsize OP},
  6/colPad!20/{\scriptsize[P]},
  7/colPad!20/{\scriptsize[P]}}{
    \node[tok, fill=\col] (tR\i) at (1.7 + \i*0.70, 6.6) {\lbl};
}
 
% ═══════════════════════════════════════════════════════════════
% RANGÉE 2 — Encodeurs
% ═══════════════════════════════════════════════════════════════
 
\node[enc, fill=colCtx!10, draw=colCtx, text=colCtx!80!black]
  (ftheta) at (-3.4, 4.5)
  {\textbf{Encodeur de Contexte}\\[2pt]$f_{\theta}$};
 
\node[enc, fill=colTgt!10, draw=colTgt, text=colTgt!80!black]
  (fxi) at (3.4, 4.5)
  {\textbf{Encodeur Cible (EMA)}\\[2pt]$f_{\xi}$};
 
% flèches séquences → encodeurs
\draw[arr, colCtx] (-3.4, 6.22) -- (ftheta.north);
\draw[arr, colTgt]  (3.4, 6.22) -- (fxi.north);
 
% ═══════════════════════════════════════════════════════════════
% FLÈCHE EMA (entre encodeurs, passant par le haut)
% ═══════════════════════════════════════════════════════════════
 
\draw[emaArr]
  (ftheta.north east) -- ++(0.5, 0.5)
  -- node[above, font=\scriptsize, colTgt!80!black]
       {$\xi \;\leftarrow\; m\,\xi + (1-m)\,\theta$}
  ++(4.2, 0) -- (fxi.north west);
 
% ═══════════════════════════════════════════════════════════════
% RANGÉE 3 — Représentations latentes
% ═══════════════════════════════════════════════════════════════
 
\node[repr, fill=colCtx!8, draw=colCtx]
  (zctx) at (-3.4, 3.0) {$\bm{Z}_{\mathrm{ctx}} \in \mathbb{R}^{L \times d}$};
 
\node[repr, fill=colTgt!8, draw=colTgt]
  (ztgt) at (3.4, 3.0)  {$\bm{Z}_{\mathrm{tgt}} \in \mathbb{R}^{L \times d}$};
 
\draw[arr, colCtx] (ftheta) -- (zctx);
\draw[arr, colTgt] (fxi)    -- (ztgt);
 
% stop-gradient sur la branche cible
\node[sgLbl] at (3.85, 3.75) {stop-grad};
\draw[colTgt!60, dashed, line width=0.5pt]
  (3.70, 3.44) -- (3.70, 4.17);
 
% ═══════════════════════════════════════════════════════════════
% RANGÉE 4 — Prédicteur (centré)
% ═══════════════════════════════════════════════════════════════
 
\node[pred, fill=colPred!12, draw=colPred, text=colPred!70!black]
  (gphi) at (0, 1.8)
  {\textbf{Prédicteur}\\[2pt]$g_{\phi}$};
 
% Z_ctx → prédicteur
\draw[arr, colCtx]
  (zctx.south) -- ++(0, -0.3) -| (gphi.north);
 
% ═══════════════════════════════════════════════════════════════
% RANGÉE 5 — Prédiction
% ═══════════════════════════════════════════════════════════════
 
\node[repr, fill=colPred!8, draw=colPred]
  (zhat) at (0, 0.5) {$\hat{\bm{Z}} \in \mathbb{R}^{L \times d}$};
 
\draw[arr, colPred] (gphi) -- (zhat);
 
% ═══════════════════════════════════════════════════════════════
% RANGÉE 6 — Perte (centrée)
% ═══════════════════════════════════════════════════════════════
 
\node[loss, fill=colLoss!8, draw=colLoss, text=colLoss!80!black]
  (loss_node) at (0, -0.85)
  {$\mathcal{L} = \mathrm{HuberLoss}\!\left(
      \hat{\bm{Z}}_{\mathcal{M}},\;
      \mathrm{sg}\!\left(\bm{Z}^{\mathrm{tgt}}_{\mathcal{M}}\right)
    \right)$};
 
% Z_hat → perte
\draw[arr, colPred] (zhat) -- (loss_node.north -| zhat);
 
% Z_tgt → perte (descend, puis va à gauche vers la perte)
\draw[arr, colTgt]
  (ztgt.south) -- ++(0,-0.35) -- ++(0,-1.35)
  -| (loss_node.east);
 
\end{tikzpicture}
\caption{Architecture I-JEPA 1D. La séquence masquée $\tilde{\bm{x}}$
  alimente l'encodeur de contexte $f_\theta$ (bleu) ; la séquence
  intacte $\bm{x}$ alimente l'encodeur cible $f_\xi$ (vert), mis à
  jour exclusivement par EMA (flèche pointillée). Le prédicteur
  $g_\phi$ (ambre) produit $\hat{\bm{Z}}$ depuis $\bm{Z}_{\text{ctx}}$,
  et la perte de Huber est calculée uniquement sur la zone cible
  $\mathcal{M}$. Aucun gradient ne se propage à travers $f_\xi$
  (\emph{stop-gradient}).}
\label{fig:ijepa_overview}
\end{figure}
 
% ============================================================
\subsubsection{Représentation d'Entrée}
\label{sub:input_repr}
% ============================================================
 
Chaque chemin extrait par le pipeline (Section~\ref{fig:pipeline_extraction})
est une liste de tokens symboliques issus de l'IR VEX et des points
de sortie vers les API système. Les tokens sont convertis en indices
entiers par le vocabulaire $\mathcal{V}$, puis les séquences sont
tronquées à $L = 500$ tokens et complétées par le token spécial
\texttt{<PAD>} (identifiant $0$) pour atteindre la longueur fixe $L$.
Formellement, pour un chemin $p$ de longueur $\ell_p$,
\[
  x_i =
  \begin{cases}
    v(p_i) & \text{si } i \le \ell_p,\\
    0       & \text{sinon (padding)},
  \end{cases}
\]
où $v : \mathcal{V} \to \{0, \ldots, |\mathcal{V}|-1\}$ est la
fonction d'encodage du vocabulaire.
 
% ============================================================
\subsubsection{Encodeur de Contexte \texorpdfstring{$f_\theta$}{f\_theta}}
\label{sub:context_encoder}
% ============================================================
 
L'encodeur de contexte projette une séquence de tokens (partiellement
masquée) dans un espace de représentation continu de dimension $d=256$.
Il est composé d'une couche d'embedding suivie d'un réseau de
convolutions 1D causales.
 
\paragraph{Embedding.}
Les indices entiers sont d'abord projetés dans $\mathbb{R}^{e}$
via une matrice apprise $\bm{E} \in \mathbb{R}^{|\mathcal{V}| \times e}$,
avec $e = 128$ :
\[
  \bm{H}^{(0)} = \bm{E}[\bm{x}] \in \mathbb{R}^{L \times e}.
\]
 
\paragraph{Réseau convolutif.}
Trois couches de convolution 1D sont appliquées successivement sur
la dimension temporelle, selon le schéma
$e \to h \to h \to d$ avec $h = 256$ et $d = 256$ :
\begin{align}
  \bm{H}^{(1)} &= \mathrm{GELU}\!\left(
    \mathrm{Conv1D}_{k=5}^{\,e \to h}\!\left(\bm{H}^{(0)}\right)\right),
  \label{eq:conv1}\\
  \bm{H}^{(2)} &= \mathrm{GELU}\!\left(
    \mathrm{Conv1D}_{k=5}^{\,h \to h}\!\left(\bm{H}^{(1)}\right)\right),
  \label{eq:conv2}\\
  \bm{Z} &= \mathrm{Conv1D}_{k=3}^{\,h \to d}\!\left(\bm{H}^{(2)}\right)
            \in \mathbb{R}^{L \times d}.
  \label{eq:conv3}
\end{align}
Les convolutions utilisent un padding ``same'' ($p = \lfloor k/2
\rfloor$), de sorte que la longueur de la séquence est préservée à
chaque étape. L'encodeur cible $f_\xi$ partage la même architecture
que $f_\theta$ mais possède ses propres paramètres, mis à jour
exclusivement par EMA (Section~\ref{sub:ema}).
 
% ------------------------------------------------------------------
% Figure 2 : architecture de l'encodeur 1D
% ------------------------------------------------------------------
\begin{figure}[ht]
\centering
\resizebox{0.82\linewidth}{!}{%
\begin{tikzpicture}[
  every node/.style={font=\small},
  layer/.style={
    draw, rounded corners=3pt,
    minimum width=1.6cm, minimum height=3.2cm,
    align=center, line width=0.8pt
  },
  arr/.style={-Stealth, line width=0.9pt},
  dim/.style={font=\tiny, fill=white, inner sep=1pt},
]
 
% embedding
\node[layer, fill=gray!10, draw=gray!60]
  (emb) at (0,0) {Embedding\\$|\mathcal{V}| \times e$};
 
% conv1
\node[layer, fill=colCtx!15, draw=colCtx, right=1.0cm of emb]
  (c1) {\texttt{Conv1D}\\$k=5,\,p=2$\\$e \to h$\\\texttt{GELU}};
 
% conv2
\node[layer, fill=colCtx!20, draw=colCtx, right=1.0cm of c1]
  (c2) {\texttt{Conv1D}\\$k=5,\,p=2$\\$h \to h$\\\texttt{GELU}};
 
% conv3
\node[layer, fill=colCtx!28, draw=colCtx, right=1.0cm of c2]
  (c3) {\texttt{Conv1D}\\$k=3,\,p=1$\\$h \to d$};
 
% sortie
\node[draw=colCtx, fill=colCtx!8, rounded corners=3pt,
      minimum width=1.6cm, minimum height=0.55cm, right=1.0cm of c3]
  (out) {$\bm{Z} \in \mathbb{R}^{L \times d}$};
 
% flèches + dimensions
\draw[arr] (emb) -- (c1)
  node[dim, midway, above] {$(B,L,e)$};
\draw[arr] (c1) -- (c2)
  node[dim, midway, above] {$(B,L,h)$};
\draw[arr] (c2) -- (c3)
  node[dim, midway, above] {$(B,L,h)$};
\draw[arr] (c3) -- (out)
  node[dim, midway, above] {$(B,L,d)$};
 
% input
\node[left=0.8cm of emb, font=\small] (inp) {$\bm{x} \in \mathbb{Z}^{B \times L}$};
\draw[arr] (inp) -- (emb);
 
% annotation hyperparamètres
\node[font=\tiny, gray, below=0.25cm of c1] {$e=128$};
\node[font=\tiny, gray, below=0.25cm of c2] {$h=256$};
\node[font=\tiny, gray, below=0.25cm of c3] {$d=256$};
 
\end{tikzpicture}
}
\caption{Architecture de l'encodeur convolutif 1D $f_\theta$
  (identique pour $f_\xi$). Les dimensions sont indiquées par batch
  $B$, longueur de séquence $L$, et canaux. Le padding \emph{same}
  préserve $L$ tout au long du réseau.}
\label{fig:encoder_arch}
\end{figure}
 
% ============================================================
\subsubsection{Prédicteur \texorpdfstring{$g_\phi$}{g\_phi}}
\label{sub:predictor}
% ============================================================
 
Le prédicteur est un réseau entièrement connecté léger opérant en
\emph{position-wise} sur les représentations produites par $f_\theta$.
Pour chaque position $i \in \{1, \ldots, L\}$ :
\[
  \hat{\bm{z}}_i
  = \bm{W}_2\,\mathrm{GELU}\!\left(\bm{W}_1\,\bm{z}^{\text{ctx}}_i
    + \bm{b}_1\right) + \bm{b}_2,
\]
avec $\bm{W}_1 \in \mathbb{R}^{d_h \times d}$,
$\bm{W}_2 \in \mathbb{R}^{d \times d_h}$, $d_h = 512$.
En notation matricielle (sur le batch entier),
\[
  \hat{\bm{Z}} = g_\phi\!\left(\bm{Z}_{\text{ctx}}\right)
  \in \mathbb{R}^{B \times L \times d}.
\]
Le prédicteur est intentionnellement peu profond : sa capacité
réduite l'empêche de mémoriser la séquence d'entrée et force
l'encodeur de contexte à produire des représentations
sémantiquement riches~\cite{assran2023self}.
 
% ============================================================
\subsubsection{Stratégie de Masquage}
\label{sub:masking}
% ============================================================
 
Contrairement aux approches de masquage uniforme (BERT, MAE), nous
adoptons un masquage par \emph{blocs contigus} tenant compte de la
longueur réelle $\ell_i$ de chaque séquence dans le batch, définie
par
\[
  \ell_i = \bigl|\{j : x_{i,j} \ne x_{\text{PAD}}\}\bigr|
         = \sum_{j=1}^{L} \mathbf{1}[x_{i,j} \ne 0].
\]
La séquence réelle est partitionnée en trois blocs de taille
approximativement égale. Deux masques booléens sont construits pour
chaque exemple $i$ du batch :
 
\begin{itemize}
  \item \textbf{Masque d'entrée} $\mathcal{M}^{\text{in}}_i$ —
    positions $[\lfloor\ell_i/3\rfloor,\, \ell_i)$ sont remplacées
    par \texttt{<MASK>} (ID~2) avant d'entrer dans $f_\theta$.
    Cela cache les blocs 2 et 3 au contexte.
  \item \textbf{Masque de prédiction} $\mathcal{M}^{\text{pred}}_i$
    — positions $[\lfloor 2\ell_i/3\rfloor,\, \ell_i)$ constituent
    la \emph{cible} : c'est sur ces positions uniquement que la
    perte est calculée.
\end{itemize}
 
Formellement, la séquence présentée à $f_\theta$ est
\[
  \tilde{x}_{i,j} =
  \begin{cases}
    x_{\text{MASK}} & \text{si } j \in \mathcal{M}^{\text{in}}_i,\\
    x_{i,j}         & \text{sinon}.
  \end{cases}
\]
La Figure~\ref{fig:masking} illustre cette stratégie en tiers sur une
séquence de longueur réelle arbitraire $\ell$.
 
% ------------------------------------------------------------------
% Figure 3 : stratégie de masquage en tiers
% ------------------------------------------------------------------
\begin{figure}[ht]
\centering
\resizebox{0.88\linewidth}{!}{%
\begin{tikzpicture}[
  every node/.style={font=\small},
  cell/.style={draw, minimum width=0.6cm, minimum height=0.6cm,
               inner sep=0pt, font=\footnotesize},
]
 
\def\L{12}   % tokens réels
\def\P{3}    % tokens padding
 
% -- Rangée 1 : séquence originale --
\node[font=\footnotesize\bfseries, anchor=east] at (-0.15, 2.0) {Séquence $\bm{x}$};
\foreach \i in {1,...,\L} {
  \node[cell, fill=colCtx!20] (o\i) at (\i*0.65-0.33, 2.0) {\i};
}
\foreach \j in {1,...,\P} {
  \node[cell, fill=colPad!20, draw=gray!40]
    at (\L*0.65+\j*0.65-0.33, 2.0) {\texttt{P}};
}
 
% -- Rangée 2 : séquence masquée (entrée f_theta) --
\node[font=\footnotesize\bfseries, anchor=east] at (-0.15, 1.1) {Entrée $\tilde{\bm{x}}$};
% tier 1 : contexte
\foreach \i in {1,...,4} {
  \node[cell, fill=colCtx!20] at (\i*0.65-0.33, 1.1) {\i};
}
% tier 2 : masqué en entrée (pré-cible)
\foreach \i in {5,...,8} {
  \node[cell, fill=colMask!40, draw=colMask] at (\i*0.65-0.33, 1.1)
    {\texttt{M}};
}
% tier 3 : masqué en entrée (cible)
\foreach \i in {9,...,12} {
  \node[cell, fill=colMask!65, draw=colMask!80] at (\i*0.65-0.33, 1.1)
    {\texttt{M}};
}
% padding
\foreach \j in {1,...,\P} {
  \node[cell, fill=colPad!20, draw=gray!40]
    at (\L*0.65+\j*0.65-0.33, 1.1) {\texttt{P}};
}
 
% -- Rangée 3 : masque de prédiction --
\node[font=\footnotesize\bfseries, anchor=east] at (-0.15, 0.2) {Cible $\mathcal{M}^{\text{pred}}$};
\foreach \i in {1,...,8} {
  \node[cell, fill=white, draw=gray!30] at (\i*0.65-0.33, 0.2) {};
}
\foreach \i in {9,...,12} {
  \node[cell, fill=colLoss!50, draw=colLoss] at (\i*0.65-0.33, 0.2)
    {\checkmark};
}
\foreach \j in {1,...,\P} {
  \node[cell, fill=colPad!20, draw=gray!40]
    at (\L*0.65+\j*0.65-0.33, 0.2) {};
}
 
% -- Accolades annotations --
\draw[decorate, decoration={brace,amplitude=5pt}]
  (0.0, 2.4) -- (4*0.65, 2.4)
  node[midway, above=5pt, font=\tiny\bfseries, colCtx!80!black]
    {Tier 1 — Contexte};
\draw[decorate, decoration={brace,amplitude=5pt}]
  (4*0.65+0.02, 2.4) -- (8*0.65, 2.4)
  node[midway, above=5pt, font=\tiny\bfseries, colMask!80!black]
    {Tier 2 — Masqué (pré-cible)};
\draw[decorate, decoration={brace,amplitude=5pt}]
  (8*0.65+0.02, 2.4) -- (12*0.65, 2.4)
  node[midway, above=5pt, font=\tiny\bfseries, colLoss!70!black]
    {Tier 3 — Cible $\mathcal{M}^{\text{pred}}$};
\draw[decorate, decoration={brace,amplitude=5pt}]
  (12*0.65+0.02, 2.4) -- (15*0.65, 2.4)
  node[midway, above=5pt, font=\tiny\bfseries, colPad!60!black]
    {Padding};
 
% -- Légende frontières --
\draw[dashed, gray] (4*0.65+0.01, -0.15) -- (4*0.65+0.01, 2.7)
  node[above, font=\tiny] {$\lfloor\ell/3\rfloor$};
\draw[dashed, gray] (8*0.65+0.01, -0.15) -- (8*0.65+0.01, 2.7)
  node[above, font=\tiny] {$\lfloor 2\ell/3\rfloor$};
\draw[dashed, gray!50] (12*0.65+0.01, -0.15) -- (12*0.65+0.01, 2.7)
  node[above, font=\tiny] {$\ell$};
 
\end{tikzpicture}
}
\caption{Stratégie de masquage par tiers. La séquence réelle de
  longueur $\ell$ est divisée en trois blocs égaux. Le Tier~2
  (pré-cible) et le Tier~3 (cible) sont remplacés par
  \texttt{<MASK>} dans l'entrée de $f_\theta$. La perte est
  calculée uniquement sur le Tier~3 ($\mathcal{M}^{\text{pred}}$).
  Le padding ne contribue ni au masquage ni à la perte.}
\label{fig:masking}
\end{figure}
 
Cette conception assure que le modèle apprend à inférer la
\emph{continuation sémantique} d'un chemin d'exécution, ce qui
est analogue à la capacité d'un analyste à anticiper quels appels
système ou opérations mémoire suivront une séquence d'instructions
observée. Elle diffère du masquage aléatoire (BERT) en forçant une
dépendance \emph{ordonnée} : le contexte est toujours antérieur à
la cible, ce qui respecte la nature directionnelle d'un chemin de
flot de contrôle.
 
% ============================================================
\subsubsection{Mise à jour par Moyenne Mobile Exponentielle}
\label{sub:ema}
% ============================================================
 
L'encodeur cible $f_\xi$ n'est jamais optimisé directement par
rétropropagation. Ses paramètres $\xi$ sont mis à jour après chaque
itération $t$ selon la règle EMA :
\begin{equation}
  \xi_t \leftarrow m \cdot \xi_{t-1} + (1-m) \cdot \theta_t,
  \label{eq:ema}
\end{equation}
où $m \in (0,1)$ est le momentum EMA, fixé à $m = 0{,}996$ dans
nos expériences. L'initialisation est $\xi_0 = \theta_0$.
 
Ce mécanisme joue un rôle fondamental dans la stabilité de
l'apprentissage : il agit comme un \emph{stop-gradient} implicite
sur la branche cible, empêchant l'effondrement de la représentation
(\emph{representation collapse}) qui surviendrait si les deux
encodeurs évoluaient librement. L'encodeur cible produit ainsi des
cibles \emph{lissées dans le temps} qui guident l'encodeur de
contexte vers des représentations cohérentes. Ce principe,
formalisé par \cite{grill2020byol} et repris dans
I-JEPA \cite{assran2023self}, est ici transposé à l'espace des
séquences symboliques.
 
La Figure~\ref{fig:ema_dynamics} illustre comment le momentum $m$
contrôle la vitesse de convergence du professeur vers l'élève.
 
% ------------------------------------------------------------------
% Figure 4 : dynamique EMA selon m
% ------------------------------------------------------------------
\begin{figure}[ht]
\centering
\begin{tikzpicture}
\begin{axis}[
  width=0.72\linewidth, height=4.8cm,
  xlabel={Itération $t$},
  ylabel={$\|\xi_t - \theta_t\|_2$ (normalisé)},
  xmin=0, xmax=300,
  ymin=0, ymax=1.05,
  xtick={0,50,100,150,200,250,300},
  ytick={0,0.25,0.5,0.75,1.0},
  grid=both, grid style={line width=0.3pt, draw=gray!25},
  axis background/.style={fill=gray!4},
  legend style={at={(0.97,0.97)}, anchor=north east,
                font=\tiny, fill=white, draw=gray!40},
  label style={font=\small},
  tick label style={font=\tiny},
]
% m = 0.90
\addplot[colMask, thick, smooth] coordinates {
  (0,1)(10,0.65)(20,0.42)(40,0.18)(60,0.08)(100,0.02)(200,0.001)(300,0)
};
\addlegendentry{$m=0{,}90$}
% m = 0.996 (notre valeur)
\addplot[colCtx, thick, smooth, dashed] coordinates {
  (0,1)(20,0.92)(50,0.82)(100,0.67)(150,0.55)(200,0.44)(300,0.29)
};
\addlegendentry{$m=0{,}996$ (utilisé)}
% m = 0.9999
\addplot[colTgt, thick, smooth, dotted] coordinates {
  (0,1)(50,0.995)(100,0.99)(200,0.98)(300,0.97)
};
\addlegendentry{$m=0{,}9999$}
\end{axis}
\end{tikzpicture}
\caption{Évolution schématique de la distance $\|\xi_t -
  \theta_t\|_2$ pour différentes valeurs de momentum $m$.
  Un $m$ trop faible (0{,}90) rend le professeur trop réactif et
  instabilise les cibles. Un $m$ trop élevé (0{,}9999) ralentit
  excessivement la convergence. La valeur $m=0{,}996$ offre le
  meilleur compromis pour notre échelle d'entraînement.}
\label{fig:ema_dynamics}
\end{figure}
 
% ============================================================
\subsubsection{Fonction de Perte}
\label{sub:loss}
% ============================================================
 
La perte est calculée exclusivement sur les positions appartenant
au masque de prédiction $\mathcal{M}^{\text{pred}}$, afin que le
modèle ne reçoive aucun signal d'apprentissage sur les positions
de contexte ou de padding. Pour un batch de taille $B$, la perte
est la \emph{Huber Loss} (ou \emph{Smooth L1 Loss}) entre les
représentations prédites et les représentations cibles :
\begin{equation}
  \mathcal{L}(\theta, \phi ; \xi) =
  \frac{1}{\sum_{i=1}^{B}|\mathcal{M}^{\text{pred}}_i|}
  \sum_{i=1}^{B}
  \sum_{j \in \mathcal{M}^{\text{pred}}_i}
  \sum_{k=1}^{d}
  \ell_{\delta}\!\left(\hat{z}_{i,j,k},\,
    \mathrm{sg}\!\left(z^{\text{tgt}}_{i,j,k}\right)\right),
  \label{eq:loss}
\end{equation}
où $\mathrm{sg}(\cdot)$ désigne l'opérateur \emph{stop-gradient}
(pas de rétropropagation à travers $f_\xi$), et $\ell_\delta$ est
la fonction de Huber de paramètre $\delta = 1$ :
\[
  \ell_\delta(a, b) =
  \begin{cases}
    \dfrac{1}{2}(a-b)^2          & \text{si } |a-b| \le \delta,\\[6pt]
    \delta\!\left(|a-b| - \dfrac{\delta}{2}\right) & \text{sinon.}
  \end{cases}
\]
Le choix de la Huber Loss par rapport à la MSE présente deux
avantages dans notre contexte : (i) elle est moins sensible aux
\emph{outliers} de représentation qui peuvent survenir sur des
binaires malformés ou fortement obfusqués, et (ii) son comportement
linéaire au-delà de $\delta$ accélère la phase initiale de
l'apprentissage lorsque les représentations sont encore éloignées.
 
Seuls les paramètres $\theta$ (encodeur de contexte) et $\phi$
(prédicteur) sont mis à jour par le gradient de $\mathcal{L}$.
Les paramètres $\xi$ de l'encodeur cible sont exclusivement mis à
jour par~\eqref{eq:ema}.
 
% ============================================================
\subsubsection{Récapitulatif et Hyperparamètres}
\label{sub:hyperparams}
% ============================================================
 
Le Tableau~\ref{tab:hyperparams} récapitule les hyperparamètres
du modèle et de l'entraînement.
 
\begin{table}[ht]
\centering
\caption{Hyperparamètres du modèle I-JEPA 1D.}
\label{tab:hyperparams}
\renewcommand{\arraystretch}{1.25}
\begin{tabular}{llr}
\toprule
\textbf{Composant} & \textbf{Paramètre} & \textbf{Valeur} \\
\midrule
\multirow{3}{*}{Encodeur ($f_\theta$, $f_\xi$)}
  & Dimension d'embedding $e$      & 128 \\
  & Dimension cachée $h$           & 256 \\
  & Dimension de sortie $d$        & 256 \\[2pt]
\multirow{2}{*}{Prédicteur ($g_\phi$)}
  & Dimension cachée $d_h$         & 512 \\
  & Activation                     & GELU \\[2pt]
\multirow{2}{*}{Masquage}
  & Stratégie                      & Tiers consécutifs \\
  & Ratio prédit / séquence        & $\lfloor\ell/3\rfloor$ \\[2pt]
\multirow{3}{*}{Entraînement}
  & Optimiseur                     & AdamW \\
  & Taux d'apprentissage $\eta$    & $10^{-4}$ \\
  & Taille de batch $B$            & 10 \\[2pt]
EMA & Momentum $m$                 & 0{,}996 \\
Perte & Fonction                   & Huber ($\delta=1$) \\
Vocabulaire & Taille $|\mathcal{V}|$ & 371 \\
\bottomrule
\end{tabular}
\end{table}
